\subsection{End Behavior} 
We can also discuss the ``end behavior'' of a polynomial. This refers to what happens to the value of $f$ when $x$ is a very large \makeblank[1.5in] number or a very large \makeblank[1.5in] number.
\begin{itemize}
\item To say, ``$x$ gets very large,'' we write \makeblank
\item To say, ``$x$ gets very large negative,'' we write \makeblank
\end{itemize}
When $x$ is very large (positive or negative), the \makeblank of the polynomial is the most important.
\begin{example}Describe the end behavior of the polynomial $$f(x)=-3x(x-5)(2x+3)(4-x)^2$$
\begin{itemize}
\item The leading term is \makeblank[4in]
\item The \makeblank[1in] is \makeblank[1in]
\item The \makeblank[1in] is \makeblank[1in]
\item The graph of $y=\makeblanksmall$ looks like this:
      \begin{ilgraph}
			\useasboundingbox(-4,-4) rectangle (4,2);
			\draw[step=1,npblue] (-4,-4) grid (4,4);
			\draw[axis] (-4.5,0) -- (4.5,0) node[anchor=south, font=\scriptsize]{$x$};
			\draw[axis] (0,-4.5) -- (0,4.5) node[anchor=west, font=\scriptsize]{$y$};
			\end{ilgraph}
\item We describe this end behavior as:
\begin{itemize} 
\item $f(x)\rightarrow\makeblanksmall$ when $x\rightarrow +\infty$
\item $f(x)\rightarrow\makeblanksmall$ when $x\rightarrow -\infty$
\end{itemize}
\end{itemize}
\end{example}
This is the easiest way to find the end behavior: find the leading term and sketch its graph.